%=============================================================================!
% 11.7  HOW FAR THESE METHODS GO                                              !
%=============================================================================!
\section{How far these methods go}
\label{sec:cs-limits}

\Cref{tab:cs-steps} collects the steps measured in this chapter on the
same 64 molecules of water, each at the energy fluctuation of 1\% of the
kinetic energy's that \cref{sec:cs-water} took as the level. The cost of
a run is set by the calls of the costly forces between molecules, one per
step, or per outer step for RESPA.

\begin{table}[htb]
\centering
\caption{The step at which the energy fluctuation of 64 water molecules
reaches 1\% of the kinetic energy's, with each method of this chapter;
the gain over flexible water by velocity Verlet; and the calls of the
forces between molecules per picosecond, from the unrounded steps. Heavy H: hydrogen three times as
heavy, at the oxygen's expense.}
\label{tab:cs-steps}
\begin{tabular}{llccc}
\toprule
model and method & fastest motion left & step / fs & gain & calls per ps\\
\midrule
flexible, velocity Verlet & O--H stretch & 0.41 & 1 & 2430\\
flexible, RESPA & O--H stretch & 0.72 & 1.8 & 1380\\
O--H held & bend & 0.97 & 2.4 & 1030\\
O--H held, heavy H & bend, slower & 1.31 & 3.2 & 765\\
rigid & turning & 1.98 & 4.8 & 505\\
rigid, heavy H & turning, slower & 3.38 & 8.2 & 296\\
\bottomrule
\end{tabular}
\end{table}

Each method removes the fastest motion or slows it, and the next fastest
then sets the step. Holding the O--H bonds leaves the bend; holding the
angle too leaves the turning of whole molecules; heavier hydrogens slow
whatever is left. RESPA keeps every motion; at the 1\% level its outer
step is \SI{0.72}{\femto\second}, and resonance holds it below about a
quarter of the fastest period, \SI{1.8}{\femto\second}, whatever level is
accepted (\cref{sec:cs-resonance}). The factors are of a few each, and they
combine: rigid water with heavier hydrogens takes steps 8.2 times as long
as the flexible model, but it is no longer the same model. For a
simulation that keeps every motion, the gain here is the 1.8 of RESPA.

A factor of thirty, which TrajCast reaches on quartz, would need steps of
\SI{12}{\femto\second} for the flexible model, longer than the period of its fastest vibration,
\SI{7.33}{\femto\second}. Such steps lie outside the range of the methods developed here.
Velocity Verlet becomes unstable beyond $2/\omega$
(\cref{sec:in-stability}), \SI{2.33}{\femto\second} here, and every
method of this chapter must still follow, with steps short compared with
its period, the fastest motion it keeps. The authors of TrajCast make the
same point about methods with multiple time steps, which are
`fundamentally limited by the need for small time steps in parts of the
system'\cite{Thiemann2026}. TrajCast itself, run on 64 molecules of
flexible water in a box of nearly the same size as this chapter's, though
with a different model of water, takes strides of $\Dt =
\SI{5}{\femto\second}$, ten times the \SI{0.5}{\femto\second} step of the
simulations it learned from, and on quartz strides of
\SI{30}{\femto\second}, thirty times the \SI{1}{\femto\second} step of its
simulations\cite{Thiemann2026}. A stride of \SI{5}{\femto\second} already
exceeds the stability limit of velocity Verlet for the faster O--H stretch
of real water, of period \SI{8.88}{\femto\second}, which is
\SI{2.83}{\femto\second} (\cref{sec:in-stability}), as well as the
\SI{2.33}{\femto\second} of this chapter's model. A learned map that predicts
the state a whole stride ahead is not bound by that limit, because it does
not step through the vibrations in between; what it must get right
instead, and how it can be trusted, is the subject of the planned Part~VI.

\begin{keyidea}
Constraints, heavier hydrogens and multiple time steps each lengthen the
step by a factor of a few, by removing or slowing the fastest motion,
until the next fastest sets the limit. For velocity Verlet, stability requires a step below $2/\omega$ for the
fastest vibration retained; this is a limit of that integrator, not of
every possible numerical method. A learned propagator is not bound by
the same limit, but its accuracy and sampled distribution still need
independent tests.
\end{keyidea}

\begin{selfcheck}
From \cref{tab:cs-steps}, how many calls of the forces between molecules
does a nanosecond of rigid water with heavier hydrogens take, and how
many does flexible water by velocity Verlet take?
\end{selfcheck}
